\documentclass[reqno,11pt]{amsart}
\usepackage{math327-lecture}
\usepackage{needspace}
\newcommand{\Q}{\mathbb{Q}}
\newcommand{\C}{\mathbb{C}}
\lectureinfo{27}{Nov. 20, 2026}{Ideals, Principal Ideals, and Quotient Rings}

% Typeset from handwritten Lecture 29, Nov. 21, 2025, pp. 1--12.
\begin{document}
\raggedbottom
\makelecturetitle

\section{Kernels of ring homomorphisms and ideals}

Throughout, rings are commutative with identity. Let
$\varphi:R\to S$ be a homomorphism of rings. Its \emph{kernel} is
\[
  \ker(\varphi)=\{x\in R:\varphi(x)=0\}.
\]

\begin{proposition}\label{prop:injective}
$\ker(\varphi)=0$ if and only if $\varphi$ is injective.
\end{proposition}
\begin{proof}
In particular, $\varphi$ is a homomorphism of the additive groups
underlying $R$ and $S$, and $\ker(\varphi)$ is its kernel as a group
homomorphism. The statement therefore follows from what we know about
kernels of group homomorphisms.
\end{proof}

The kernel $\ker(\varphi)\subseteq R$ is not just an additive
subgroup: it is an \emph{ideal} of $R$.

\begin{define}[Ideal]
A subset $I\subseteq R$ is an \emph{ideal} if
\begin{enumerate}[label=(\arabic*)]
  \item $0\in I$;
  \item if $a,b\in I$, then $a+b\in I$;
  \item if $a\in I$ and $b\in R$, then $ab\in I$.
\end{enumerate}
\end{define}

Note that in condition (3), the element $b$ is allowed to be any
element of $R$, not just of $I$. For the kernel, condition (3) holds
because $\varphi(ab)=\varphi(a)\varphi(b)=0\cdot\varphi(b)=0$
whenever $a\in\ker(\varphi)$.

\begin{example}[Ideals of $\Z$]
In the ring $\Z$, an ideal $I$ could be
\[
  \langle n\rangle=\{0,\pm n,\pm2n,\dots\}=n\Z,
\]
and in fact these are all the ideals of $\Z$.
\end{example}

\section{Ideals generated by elements}

\begin{definition}[Ideal generated by elements]
If $f_1,\dots,f_n\in R$, we write $I=(f_1,\dots,f_n)$ for the
\emph{ideal generated by $f_1,\dots,f_n$}, that is, the smallest
ideal of $R$ containing $f_1,\dots,f_n$.
\end{definition}

What is it? It is the set of all combinations
\[
  (f_1,\dots,f_n)=\{a_1f_1+\cdots+a_nf_n:a_i\in R\}.
\]
Any ideal containing $f_1,\dots,f_n$ contains all such combinations
by conditions (2) and (3), so it remains to check that this set $I$
is an ideal. Clearly $0\in I$, and $I$ is closed under addition.
Moreover, if $b\in R$, then
\[
  b(a_1f_1+\cdots+a_nf_n)=(ba_1)f_1+\cdots+(ba_n)f_n\in I.
\]

\begin{example}
\begin{enumerate}[label=(\alph*)]
  \item $(1)=R$.
  \item For $a\in R$, we have $(a)=\{ba:b\in R\}=Ra$.
\end{enumerate}
\end{example}

Recall from last time that
\[
  R\text{ is a field}\iff (0)\text{ and }(1)\text{ are the only ideals of }R.
\]

\section{Every ideal of \texorpdfstring{$F[x]$}{F[x]} is principal}

Let $F$ be a field (so, for instance, $F=\Q$, $\R$, $\C$,
$\Z/p\Z$, \dots), and let $R=F[x]$ be the ring of polynomials with
coefficients in $F$. What are the ideals of $R$? One example is
$(f(x))$ for any $f\in F[x]$.

\begin{define}[Principal ideal]
If an ideal has the form $I=(f_1,\dots,f_n)$, we say that $I$ is
\emph{generated by} $f_1,\dots,f_n$. If $I=(f)$ for a single element
$f$, we say that $I$ is \emph{principal}.
\end{define}

\begin{theorem}\label{thm:Fx-principal}
Every ideal in $F[x]$ is principal.
\end{theorem}

This is another instance of the analogy $F[x]\longleftrightarrow\Z$.

\begin{proof}
Let $I\subseteq F[x]$ be any ideal. If $I=(0)$, it is principal.
If $I\neq(0)$, then $I$ contains a nonzero polynomial of smallest
degree; call it $f\in I$, say
\[
  f(x)=a_nx^n+a_{n-1}x^{n-1}+\cdots+a_0,\qquad a_n\neq0.
\]
Since $F$ is a field and $f\in I$, so is $\frac1{a_n}f\in I$.
That is, there exists a monic polynomial
\[
  x^n+b_{n-1}x^{n-1}+\cdots+b_0\in I
\]
of smallest degree; replace $f$ by it.

Now let $g(x)\in I$. By the division algorithm for $F[x]$, since $f$
is monic, we may write
\[
  g(x)=h(x)f(x)+r(x)
\]
for some $r(x)$ which is either $0$ or of degree less than
$\deg(f)$. Since $g\in I$ and $hf\in I$,
\[
  g-hf=r\in I.
\]
If $r\neq0$, then $\deg(r)<\deg(f)$, contradicting the choice of $f$
as an element of smallest degree in $I$. Therefore $r=0$ and
$g=h\cdot f$, i.e.\ $g\in(f)$. Hence $I=(f)$.
\end{proof}

\begin{warningbox}
If $R$ is a general ring (not a field), then $R[x]$ need not be a
principal ideal domain; that is, not every ideal of $R[x]$ need be
principal.
\end{warningbox}

\begin{example}[A non-principal ideal of {$\Z[x]$}]
In $\Z[x]$, consider $I=(3,5x)$. Can we generate $I$ with just one
element? Suppose $I=(g)$ with $g\in\Z[x]$. Since $3\in(g)$, we have
$3=hg$ for some $h$, so $g$ must have degree $0$, i.e.\ $g$ is a
constant. Also $g\mid3$ and $g\mid5x$. The constant divisors of $3$
are $\pm1,\pm3$, and $\pm3$ does not divide $5x$. So $g=\pm1$ and
$I=\Z[x]$. But then
\[
  1=a(x)\cdot3+b(x)\cdot5x
\]
for some $a,b\in\Z[x]$, which is impossible: comparing constant
terms gives $1=3a(0)$ in $\Z$.
% NOTE: the notes say "impossible for degree reasons"; the argument
% given here compares constant terms instead.
Hence $I$ is not principal.
\end{example}

\section{A kernel in \texorpdfstring{$\Q[x]$}{Q[x]}}

By Theorem~\ref{thm:Fx-principal}, $\Q[x]$ has only principal ideals.

\begin{example}[The kernel of evaluation at {$\sqrt[3]{2}$}]
Let $\gamma=\sqrt[3]{2}$ be the real cube root of $2$, so
$\gamma^3=2$, and consider the evaluation homomorphism
\[
  \varphi:\Q[x]\to\C,\qquad x\mapsto\gamma.
\]
What is $\ker(\varphi)$? Is it $0$? We know $\ker(\varphi)=(k)$ for
some $k\in\Q[x]$. First, $x^3-2\in\ker(\varphi)$, since
\[
  \varphi(x^3-2)=\gamma^3-2=0.
\]
If $x^3-2\in\ker(\varphi)$, then $g\cdot(x^3-2)\in\ker(\varphi)$
for every $g\in\Q[x]$, so $(x^3-2)\subseteq\ker(\varphi)$. But is
there any element in $\ker(\varphi)$ of lower positive degree?

Just as in the correspondence $F[x]\longleftrightarrow\Z$, where
\[
  n\mid m\iff(m)\subseteq(n),
\]
we have for polynomials
\[
  f\mid g\iff(g)\subseteq(f).
\]
So
\[
  (x^3-2)\subseteq\ker(\varphi)=(k)\iff k\mid x^3-2,\quad k\in\Q[x].
\]
The polynomial $k$ is not a nonzero constant, since then
$\ker(\varphi)=\Q[x]$, whereas $\varphi(1)=1\neq0$. If $k$ had
degree $1$ or $2$, we could factor
\[
  x^3-2=(a_0+a_1x)(b_0+b_1x+b_2x^2),\qquad a_i,b_j\in\Q.
\]
This is impossible to solve with rational numbers
(Exercise~\ref{ex:cube-root}). Hence $\deg k=3$, so $k$ is a
nonzero scalar multiple of $x^3-2$ and
\[
  \ker(\varphi)=(x^3-2).
\]
\end{example}

\begin{exercise}\label{ex:cube-root}
Show that $x^3-2$ cannot be written as
$(a_0+a_1x)(b_0+b_1x+b_2x^2)$ with $a_i,b_j\in\Q$.
\end{exercise}

\section{Quotient rings}

Let $R$ be a ring and $I\subseteq R$ an ideal. Then we can form a
new ring $R/I$, which consists of the cosets
\[
  a+I,\qquad a\in R.
\]
As an additive group, this is the quotient group of $R$ by the
subgroup $I$. What is the multiplication? We define
\[
  (a+I)\cdot(b+I)\coloneqq ab+I.
\]

\begin{proposition}\label{prop:well-defined}
This gives a well-defined multiplication on $R/I$.
\end{proposition}
\begin{proof}
Suppose $a+I=\tilde a+I$ and $b+I=\tilde b+I$, i.e.\
$a-\tilde a\in I$ and $b-\tilde b\in I$. Then
\[
  (a+I)(b+I)=ab+I,\qquad(\tilde a+I)(\tilde b+I)=\tilde a\tilde b+I,
\]
and we must show $ab+I=\tilde a\tilde b+I$, i.e.\
$ab-\tilde a\tilde b\in I$. Why does this hold? We have
\[
  \tilde b(a-\tilde a)+a(b-\tilde b)=ab-\tilde a\tilde b,
\]
and both terms on the left lie in $I$ by condition (3) in the
definition of an ideal. Since $I$ is closed under addition,
$ab-\tilde a\tilde b\in I$.
\end{proof}

\begin{exercise}
Check that $R/I$ is a ring with
\begin{align*}
  \text{addition: }&(a+I)+(b+I)=(a+b)+I,\\
  \text{multiplication: }&(a+I)(b+I)=ab+I,
\end{align*}
zero element $0+I$, and identity element $1+I$.
\end{exercise}

\begin{define}[Quotient ring]
The ring $R/I$ with the operations above is the \emph{quotient ring}
of $R$ by the ideal $I$.
\end{define}

\begin{example}[{$\Z/n\Z$}]
Take $R=\Z$ and $I=(n)$. Then $R/I=\Z/n\Z$, with multiplication of
residue classes $\bar a\cdot\bar b=\overline{ab}$. For instance, when
$n=12$,
\[
  \bar5\cdot\bar6=\overline{30}=\bar6.
\]
\end{example}

\end{document}
